\section{天启至崇祯的补充编年节点（一）}
这一组条目把此前尚未设导航入口的年度材料接回本章脉络。它们不是把同年材料机械等同为因果，而是在可核查的时间位置上，分别保留行动、行政记录与社会经验之间的距离。

\subsection{1626年前后的材料线索}

\eventanchor{ML-1626-006}{明·1626（天启六年）}
\paragraph{1626（天启六年）：旱灾、蝗灾、疫病、河患与赈济的本年材料连接环境和社会压力，死亡与流离数字须谨慎处理。（材料核查重点：报称信息的来源、抄传与行政分类）} 这一节点应放回本章已经展开的权力、财政、地方社会与边地关系中理解。账本把它出现的条件概括为：自然风险与地方报告促成朝廷或地方的应对记录。 随后可追踪的变化是：赈济、迁徙和损失的实际范围仍须以受灾地区材料分别核验。 可由《熹宗实录》天启六年条；《明史·五行志》；受灾府县地方志与奏疏 互相核对；不过，桥接条目只标示可回查的年度问题与材料链；实录有中央选择性，崇祯期尤其需要奏疏、地方、朝鲜及满文材料交叉。；本条特别核对报称信息的来源、抄传与行政分类。 因而这里标示的是材料能够支持的行动与制度位置，而不是对动机、地方执行或普通人经验的无保留复原。

\eventanchor{ML-1626-007}{明·1626（天启六年）}
\paragraph{1626（天启六年）：辽饷、剿饷、练饷及地方征解的本年文书显示财政负担，定额、征收与实际流向不得混为一谈。（材料核查重点：报称信息的来源、抄传与行政分类）} 这一节点应放回本章已经展开的权力、财政、地方社会与边地关系中理解。账本把它出现的条件概括为：财政征解、交通工程或货币支付的压力使相关措施进入政务。 随后可追踪的变化是：其法定安排不等于地方执行，后续须以地域材料追踪负担与成效。 可由《熹宗实录》天启六年条；《明会典》相关条目；地方志、黄册/契约/河工材料 互相核对；不过，桥接条目只标示可回查的年度问题与材料链；实录有中央选择性，崇祯期尤其需要奏疏、地方、朝鲜及满文材料交叉。；本条特别核对报称信息的来源、抄传与行政分类。 因而这里标示的是材料能够支持的行动与制度位置，而不是对动机、地方执行或普通人经验的无保留复原。

\eventanchor{ML-1626-008}{明·1626（天启六年）}
\paragraph{1626（天启六年）：朝鲜、辽东和海上关系的本年材料提供外部观察位置，明、后金（清）与地方势力的称谓需具体化。（材料核查重点：报称信息的来源、抄传与行政分类）} 这一节点应放回本章已经展开的权力、财政、地方社会与边地关系中理解。账本把它出现的条件概括为：朝贡名义、边境安全与港口贸易的利益使交涉持续发生。 随后可追踪的变化是：它为后续册封、互市或海防安排留下文书与跨政权比较线索。 可由《熹宗实录》天启六年条；《明史·外国传》；《朝鲜王朝实录》、琉球《历代宝案》或海外档案 互相核对；不过，桥接条目只标示可回查的年度问题与材料链；实录有中央选择性，崇祯期尤其需要奏疏、地方、朝鲜及满文材料交叉。；本条特别核对报称信息的来源、抄传与行政分类。 因而这里标示的是材料能够支持的行动与制度位置，而不是对动机、地方执行或普通人经验的无保留复原。

\eventanchor{ML-1627-001}{明·1627（天启七年）春}
\paragraph{1627（天启七年）春：宁晋之战：皇太极率领的后金军队进攻锦州，但被击退} 这一节点应放回本章已经展开的权力、财政、地方社会与边地关系中理解。账本把它出现的条件概括为：前期边防、地方武装或跨区域资源调动的压力推动了该行动。 随后可追踪的变化是：它改变了后续调兵、补给或地方秩序的条件；战果和伤亡须分开核对。 可由《熹宗实录》天启七年条；《明史·兵志》及相关列传；边镇、地方志或对方材料 互相核对；不过，译写候选以编年材料定位；实录记载反映朝廷选择，崇祯期须另以奏疏、地方及敌对政权材料交叉。 因而这里标示的是材料能够支持的行动与制度位置，而不是对动机、地方执行或普通人经验的无保留复原。

\eventanchor{ML-1627-002}{明·1627（天启七年）九月30日}
\paragraph{1627（天启七年）九月30日：天启皇帝驾崩} 这一节点应放回本章已经展开的权力、财政、地方社会与边地关系中理解。账本把它出现的条件概括为：继承、官僚协调或皇权运作的压力使问题进入朝廷程序。 随后可追踪的变化是：它影响后续人事与政策议程；动机和派系标签不能由单一叙事裁定。 可由《熹宗实录》天启七年条；《明史·本纪》及相关列传；诏令、奏疏或会典版本 互相核对；不过，译写候选以编年材料定位；实录记载反映朝廷选择，崇祯期须另以奏疏、地方及敌对政权材料交叉。 因而这里标示的是材料能够支持的行动与制度位置，而不是对动机、地方执行或普通人经验的无保留复原。

\eventanchor{ML-1627-003}{明·1627（天启七年）十月2日}
\paragraph{1627（天启七年）十月2日：朱由检即位崇祯皇帝} 这一节点应放回本章已经展开的权力、财政、地方社会与边地关系中理解。账本把它出现的条件概括为：继承、官僚协调或皇权运作的压力使问题进入朝廷程序。 随后可追踪的变化是：它影响后续人事与政策议程；动机和派系标签不能由单一叙事裁定。 可由《熹宗实录》天启七年条；《明史·本纪》及相关列传；诏令、奏疏或会典版本 互相核对；不过，译写候选以编年材料定位；实录记载反映朝廷选择，崇祯期须另以奏疏、地方及敌对政权材料交叉。 因而这里标示的是材料能够支持的行动与制度位置，而不是对动机、地方执行或普通人经验的无保留复原。

\eventanchor{ML-1627-004}{明·1627（天启七年）}
\paragraph{1627（天启七年）：辽东防务、后金（清）军事行动和流动作战集团的本年记录标记多战场互动，军数和胜败须保留分歧。（材料核查重点：诏令或奏议的形成与发受链）} 这一节点应放回本章已经展开的权力、财政、地方社会与边地关系中理解。账本把它出现的条件概括为：前期边防、地方武装或跨区域资源调动的压力推动了该行动。 随后可追踪的变化是：它改变了后续调兵、补给或地方秩序的条件；战果和伤亡须分开核对。 可由《熹宗实录》天启七年条；《明史·兵志》及相关列传；边镇、地方志或对方材料 互相核对；不过，桥接条目只标示可回查的年度问题与材料链；实录有中央选择性，崇祯期尤其需要奏疏、地方、朝鲜及满文材料交叉。；本条特别核对诏令或奏议的形成与发受链。 因而这里标示的是材料能够支持的行动与制度位置，而不是对动机、地方执行或普通人经验的无保留复原。

\eventanchor{ML-1627-005}{明·1627（天启七年）}
\paragraph{1627（天启七年）：天启末至崇祯朝的任免、奏疏和廷议构成本年中枢危机线索；崇祯无实录，必须以多类材料互校。（材料核查重点：诏令或奏议的形成与发受链）} 这一节点应放回本章已经展开的权力、财政、地方社会与边地关系中理解。账本把它出现的条件概括为：继承、官僚协调或皇权运作的压力使问题进入朝廷程序。 随后可追踪的变化是：它影响后续人事与政策议程；动机和派系标签不能由单一叙事裁定。 可由《熹宗实录》天启七年条；《明史·本纪》及相关列传；诏令、奏疏或会典版本 互相核对；不过，桥接条目只标示可回查的年度问题与材料链；实录有中央选择性，崇祯期尤其需要奏疏、地方、朝鲜及满文材料交叉。；本条特别核对诏令或奏议的形成与发受链。 因而这里标示的是材料能够支持的行动与制度位置，而不是对动机、地方执行或普通人经验的无保留复原。

\eventanchor{ML-1627-006}{明·1627（天启七年）}
\paragraph{1627（天启七年）：天启末至崇祯朝的任免、奏疏和廷议构成本年中枢危机线索；崇祯无实录，必须以多类材料互校。（材料核查重点：报称信息的来源、抄传与行政分类）} 这一节点应放回本章已经展开的权力、财政、地方社会与边地关系中理解。账本把它出现的条件概括为：继承、官僚协调或皇权运作的压力使问题进入朝廷程序。 随后可追踪的变化是：它影响后续人事与政策议程；动机和派系标签不能由单一叙事裁定。 可由《熹宗实录》天启七年条；《明史·本纪》及相关列传；诏令、奏疏或会典版本 互相核对；不过，桥接条目只标示可回查的年度问题与材料链；实录有中央选择性，崇祯期尤其需要奏疏、地方、朝鲜及满文材料交叉。；本条特别核对报称信息的来源、抄传与行政分类。 因而这里标示的是材料能够支持的行动与制度位置，而不是对动机、地方执行或普通人经验的无保留复原。

\eventanchor{ML-1627-007}{明·1627（天启七年）}
\paragraph{1627（天启七年）：辽东防务、后金（清）军事行动和流动作战集团的本年记录标记多战场互动，军数和胜败须保留分歧。（材料核查重点：报称信息的来源、抄传与行政分类）} 这一节点应放回本章已经展开的权力、财政、地方社会与边地关系中理解。账本把它出现的条件概括为：前期边防、地方武装或跨区域资源调动的压力推动了该行动。 随后可追踪的变化是：它改变了后续调兵、补给或地方秩序的条件；战果和伤亡须分开核对。 可由《熹宗实录》天启七年条；《明史·兵志》及相关列传；边镇、地方志或对方材料 互相核对；不过，桥接条目只标示可回查的年度问题与材料链；实录有中央选择性，崇祯期尤其需要奏疏、地方、朝鲜及满文材料交叉。；本条特别核对报称信息的来源、抄传与行政分类。 因而这里标示的是材料能够支持的行动与制度位置，而不是对动机、地方执行或普通人经验的无保留复原。

\eventanchor{ML-1627-008}{明·1627（天启七年）}
\paragraph{1627（天启七年）：地方结社、宗教、出版或士人文集的本年线索提供战乱中的社会观察，幸存者记忆不能概括全国。} 这一节点应放回本章已经展开的权力、财政、地方社会与边地关系中理解。账本把它出现的条件概括为：制度、教育或知识传播的需要推动了相关行动或文本形成。 随后可追踪的变化是：它提供制度和传播史的锚点，不能据此推断社会整体接受。 可由《熹宗实录》天启七年条；《明史·选举志》或《艺文志》；文集、碑刻或书目材料 互相核对；不过，桥接条目只标示可回查的年度问题与材料链；实录有中央选择性，崇祯期尤其需要奏疏、地方、朝鲜及满文材料交叉。 因而这里标示的是材料能够支持的行动与制度位置，而不是对动机、地方执行或普通人经验的无保留复原。

\eventanchor{ML-1628-001}{明·1628（崇祯元年）春}
\paragraph{1628（崇祯元年）春：山西遭遇干旱} 这一节点应放回本章已经展开的权力、财政、地方社会与边地关系中理解。账本把它出现的条件概括为：自然风险与地方报告促成朝廷或地方的应对记录。 随后可追踪的变化是：赈济、迁徙和损失的实际范围仍须以受灾地区材料分别核验。 可由《崇祯长编》崇祯元年条及《明史·思宗本纪》；《明史·五行志》；受灾府县地方志与奏疏 互相核对；不过，译写候选以编年材料定位；实录记载反映朝廷选择，崇祯期须另以奏疏、地方及敌对政权材料交叉。 因而这里标示的是材料能够支持的行动与制度位置，而不是对动机、地方执行或普通人经验的无保留复原。

\eventanchor{ML-1628-002}{明·1628（崇祯元年）八月}
\paragraph{1628（崇祯元年）八月：海盗头子郑芝龙投降明朝} 这一节点应放回本章已经展开的权力、财政、地方社会与边地关系中理解。账本把它出现的条件概括为：前期边防、地方武装或跨区域资源调动的压力推动了该行动。 随后可追踪的变化是：它改变了后续调兵、补给或地方秩序的条件；战果和伤亡须分开核对。 可由《崇祯长编》崇祯元年条及《明史·思宗本纪》；《明史·兵志》及相关列传；边镇、地方志或对方材料 互相核对；不过，译写候选以编年材料定位；实录记载反映朝廷选择，崇祯期须另以奏疏、地方及敌对政权材料交叉。 因而这里标示的是材料能够支持的行动与制度位置，而不是对动机、地方执行或普通人经验的无保留复原。

\eventanchor{ML-1628-003}{明·1628（崇祯元年）}
\paragraph{1628（崇祯元年）：地方结社、宗教、出版或士人文集的本年线索提供战乱中的社会观察，幸存者记忆不能概括全国。} 这一节点应放回本章已经展开的权力、财政、地方社会与边地关系中理解。账本把它出现的条件概括为：制度、教育或知识传播的需要推动了相关行动或文本形成。 随后可追踪的变化是：它提供制度和传播史的锚点，不能据此推断社会整体接受。 可由《崇祯长编》崇祯元年条及《明史·思宗本纪》；《明史·选举志》或《艺文志》；文集、碑刻或书目材料 互相核对；不过，桥接条目只标示可回查的年度问题与材料链；实录有中央选择性，崇祯期尤其需要奏疏、地方、朝鲜及满文材料交叉。 因而这里标示的是材料能够支持的行动与制度位置，而不是对动机、地方执行或普通人经验的无保留复原。

\eventanchor{ML-1638-004}{明·1638（崇祯十一年）}
\paragraph{1638（崇祯十一年）：旱灾、蝗灾、疫病、河患与赈济的本年材料连接环境和社会压力，死亡与流离数字须谨慎处理。（材料核查重点：诏令或奏议的形成与发受链）} 这一节点应放回本章已经展开的权力、财政、地方社会与边地关系中理解。账本把它出现的条件概括为：自然风险与地方报告促成朝廷或地方的应对记录。 随后可追踪的变化是：赈济、迁徙和损失的实际范围仍须以受灾地区材料分别核验。 可由《崇祯长编》崇祯十一年条及《明史·思宗本纪》；《明史·五行志》；受灾府县地方志与奏疏 互相核对；不过，桥接条目只标示可回查的年度问题与材料链；实录有中央选择性，崇祯期尤其需要奏疏、地方、朝鲜及满文材料交叉。；本条特别核对诏令或奏议的形成与发受链。 因而这里标示的是材料能够支持的行动与制度位置，而不是对动机、地方执行或普通人经验的无保留复原。

\eventanchor{ML-1638-005}{明·1638（崇祯十一年）}
\paragraph{1638（崇祯十一年）：地方结社、宗教、出版或士人文集的本年线索提供战乱中的社会观察，幸存者记忆不能概括全国。（材料核查重点：报称信息的来源、抄传与行政分类）} 这一节点应放回本章已经展开的权力、财政、地方社会与边地关系中理解。账本把它出现的条件概括为：制度、教育或知识传播的需要推动了相关行动或文本形成。 随后可追踪的变化是：它提供制度和传播史的锚点，不能据此推断社会整体接受。 可由《崇祯长编》崇祯十一年条及《明史·思宗本纪》；《明史·选举志》或《艺文志》；文集、碑刻或书目材料 互相核对；不过，桥接条目只标示可回查的年度问题与材料链；实录有中央选择性，崇祯期尤其需要奏疏、地方、朝鲜及满文材料交叉。；本条特别核对报称信息的来源、抄传与行政分类。 因而这里标示的是材料能够支持的行动与制度位置，而不是对动机、地方执行或普通人经验的无保留复原。

\eventanchor{ML-1638-006}{明·1638（崇祯十一年）}
\paragraph{1638（崇祯十一年）：旱灾、蝗灾、疫病、河患与赈济的本年材料连接环境和社会压力，死亡与流离数字须谨慎处理。（材料核查重点：报称信息的来源、抄传与行政分类）} 这一节点应放回本章已经展开的权力、财政、地方社会与边地关系中理解。账本把它出现的条件概括为：自然风险与地方报告促成朝廷或地方的应对记录。 随后可追踪的变化是：赈济、迁徙和损失的实际范围仍须以受灾地区材料分别核验。 可由《崇祯长编》崇祯十一年条及《明史·思宗本纪》；《明史·五行志》；受灾府县地方志与奏疏 互相核对；不过，桥接条目只标示可回查的年度问题与材料链；实录有中央选择性，崇祯期尤其需要奏疏、地方、朝鲜及满文材料交叉。；本条特别核对报称信息的来源、抄传与行政分类。 因而这里标示的是材料能够支持的行动与制度位置，而不是对动机、地方执行或普通人经验的无保留复原。

\eventanchor{ML-1638-007}{明·1638（崇祯十一年）}
\paragraph{1638（崇祯十一年）：辽饷、剿饷、练饷及地方征解的本年文书显示财政负担，定额、征收与实际流向不得混为一谈。} 这一节点应放回本章已经展开的权力、财政、地方社会与边地关系中理解。账本把它出现的条件概括为：财政征解、交通工程或货币支付的压力使相关措施进入政务。 随后可追踪的变化是：其法定安排不等于地方执行，后续须以地域材料追踪负担与成效。 可由《崇祯长编》崇祯十一年条及《明史·思宗本纪》；《明会典》相关条目；地方志、黄册/契约/河工材料 互相核对；不过，桥接条目只标示可回查的年度问题与材料链；实录有中央选择性，崇祯期尤其需要奏疏、地方、朝鲜及满文材料交叉。 因而这里标示的是材料能够支持的行动与制度位置，而不是对动机、地方执行或普通人经验的无保留复原。

\eventanchor{ML-1639-001}{明·1639（崇祯十二年）}
\paragraph{1639（崇祯十二年）：西班牙人和菲律宾人在吕宋岛屠杀了两万名中国人} 这一节点应放回本章已经展开的权力、财政、地方社会与边地关系中理解。账本把它出现的条件概括为：朝贡名义、边境安全与港口贸易的利益使交涉持续发生。 随后可追踪的变化是：它为后续册封、互市或海防安排留下文书与跨政权比较线索。 可由《崇祯长编》崇祯十二年条及《明史·思宗本纪》；《明史·外国传》；《朝鲜王朝实录》、琉球《历代宝案》或海外档案 互相核对；不过，译写候选以编年材料定位；实录记载反映朝廷选择，崇祯期须另以奏疏、地方及敌对政权材料交叉。 因而这里标示的是材料能够支持的行动与制度位置，而不是对动机、地方执行或普通人经验的无保留复原。

\eventanchor{ML-1639-002}{明·1639（崇祯十二年）}
\paragraph{1639（崇祯十二年）：禁止来自澳门的葡萄牙商人进入长崎} 这一节点应放回本章已经展开的权力、财政、地方社会与边地关系中理解。账本把它出现的条件概括为：朝贡名义、边境安全与港口贸易的利益使交涉持续发生。 随后可追踪的变化是：它为后续册封、互市或海防安排留下文书与跨政权比较线索。 可由《崇祯长编》崇祯十二年条及《明史·思宗本纪》；《明史·外国传》；《朝鲜王朝实录》、琉球《历代宝案》或海外档案 互相核对；不过，译写候选以编年材料定位；实录记载反映朝廷选择，崇祯期须另以奏疏、地方及敌对政权材料交叉。 因而这里标示的是材料能够支持的行动与制度位置，而不是对动机、地方执行或普通人经验的无保留复原。

\subsection{1639年前后的材料线索}

\eventanchor{ML-1639-003}{明·1639（崇祯十二年）}
\paragraph{1639（崇祯十二年）：浙江遭遇旱灾} 这一节点应放回本章已经展开的权力、财政、地方社会与边地关系中理解。账本把它出现的条件概括为：自然风险与地方报告促成朝廷或地方的应对记录。 随后可追踪的变化是：赈济、迁徙和损失的实际范围仍须以受灾地区材料分别核验。 可由《崇祯长编》崇祯十二年条及《明史·思宗本纪》；《明史·五行志》；受灾府县地方志与奏疏 互相核对；不过，译写候选以编年材料定位；实录记载反映朝廷选择，崇祯期须另以奏疏、地方及敌对政权材料交叉。 因而这里标示的是材料能够支持的行动与制度位置，而不是对动机、地方执行或普通人经验的无保留复原。

\eventanchor{ML-1639-004}{明·1639（崇祯十二年）}
\paragraph{1639（崇祯十二年）：天启末至崇祯朝的任免、奏疏和廷议构成本年中枢危机线索；崇祯无实录，必须以多类材料互校。（材料核查重点：诏令或奏议的形成与发受链）} 这一节点应放回本章已经展开的权力、财政、地方社会与边地关系中理解。账本把它出现的条件概括为：继承、官僚协调或皇权运作的压力使问题进入朝廷程序。 随后可追踪的变化是：它影响后续人事与政策议程；动机和派系标签不能由单一叙事裁定。 可由《崇祯长编》崇祯十二年条及《明史·思宗本纪》；《明史·本纪》及相关列传；诏令、奏疏或会典版本 互相核对；不过，桥接条目只标示可回查的年度问题与材料链；实录有中央选择性，崇祯期尤其需要奏疏、地方、朝鲜及满文材料交叉。；本条特别核对诏令或奏议的形成与发受链。 因而这里标示的是材料能够支持的行动与制度位置，而不是对动机、地方执行或普通人经验的无保留复原。

\eventanchor{ML-1639-005}{明·1639（崇祯十二年）}
\paragraph{1639（崇祯十二年）：旱灾、蝗灾、疫病、河患与赈济的本年材料连接环境和社会压力，死亡与流离数字须谨慎处理。} 这一节点应放回本章已经展开的权力、财政、地方社会与边地关系中理解。账本把它出现的条件概括为：自然风险与地方报告促成朝廷或地方的应对记录。 随后可追踪的变化是：赈济、迁徙和损失的实际范围仍须以受灾地区材料分别核验。 可由《崇祯长编》崇祯十二年条及《明史·思宗本纪》；《明史·五行志》；受灾府县地方志与奏疏 互相核对；不过，桥接条目只标示可回查的年度问题与材料链；实录有中央选择性，崇祯期尤其需要奏疏、地方、朝鲜及满文材料交叉。 因而这里标示的是材料能够支持的行动与制度位置，而不是对动机、地方执行或普通人经验的无保留复原。

\eventanchor{ML-1639-006}{明·1639（崇祯十二年）}
\paragraph{1639（崇祯十二年）：天启末至崇祯朝的任免、奏疏和廷议构成本年中枢危机线索；崇祯无实录，必须以多类材料互校。（材料核查重点：报称信息的来源、抄传与行政分类）} 这一节点应放回本章已经展开的权力、财政、地方社会与边地关系中理解。账本把它出现的条件概括为：继承、官僚协调或皇权运作的压力使问题进入朝廷程序。 随后可追踪的变化是：它影响后续人事与政策议程；动机和派系标签不能由单一叙事裁定。 可由《崇祯长编》崇祯十二年条及《明史·思宗本纪》；《明史·本纪》及相关列传；诏令、奏疏或会典版本 互相核对；不过，桥接条目只标示可回查的年度问题与材料链；实录有中央选择性，崇祯期尤其需要奏疏、地方、朝鲜及满文材料交叉。；本条特别核对报称信息的来源、抄传与行政分类。 因而这里标示的是材料能够支持的行动与制度位置，而不是对动机、地方执行或普通人经验的无保留复原。

\eventanchor{ML-1639-007}{明·1639（崇祯十二年）}
\paragraph{1639（崇祯十二年）：辽东防务、后金（清）军事行动和流动作战集团的本年记录标记多战场互动，军数和胜败须保留分歧。} 这一节点应放回本章已经展开的权力、财政、地方社会与边地关系中理解。账本把它出现的条件概括为：前期边防、地方武装或跨区域资源调动的压力推动了该行动。 随后可追踪的变化是：它改变了后续调兵、补给或地方秩序的条件；战果和伤亡须分开核对。 可由《崇祯长编》崇祯十二年条及《明史·思宗本纪》；《明史·兵志》及相关列传；边镇、地方志或对方材料 互相核对；不过，桥接条目只标示可回查的年度问题与材料链；实录有中央选择性，崇祯期尤其需要奏疏、地方、朝鲜及满文材料交叉。 因而这里标示的是材料能够支持的行动与制度位置，而不是对动机、地方执行或普通人经验的无保留复原。

\eventanchor{ML-1640-001}{明·1640（崇祯十三年）夏}
\paragraph{1640（崇祯十三年）夏：叛军进入四川} 这一节点应放回本章已经展开的权力、财政、地方社会与边地关系中理解。账本把它出现的条件概括为：继承、官僚协调或皇权运作的压力使问题进入朝廷程序。 随后可追踪的变化是：它影响后续人事与政策议程；动机和派系标签不能由单一叙事裁定。 可由《崇祯长编》崇祯十三年条及《明史·思宗本纪》；《明史·本纪》及相关列传；诏令、奏疏或会典版本 互相核对；不过，译写候选以编年材料定位；实录记载反映朝廷选择，崇祯期须另以奏疏、地方及敌对政权材料交叉。 因而这里标示的是材料能够支持的行动与制度位置，而不是对动机、地方执行或普通人经验的无保留复原。

\eventanchor{ML-1640-002}{明·1640（崇祯十三年）}
\paragraph{1640（崇祯十三年）：辽饷、剿饷、练饷及地方征解的本年文书显示财政负担，定额、征收与实际流向不得混为一谈。（材料核查重点：诏令或奏议的形成与发受链）} 这一节点应放回本章已经展开的权力、财政、地方社会与边地关系中理解。账本把它出现的条件概括为：财政征解、交通工程或货币支付的压力使相关措施进入政务。 随后可追踪的变化是：其法定安排不等于地方执行，后续须以地域材料追踪负担与成效。 可由《崇祯长编》崇祯十三年条及《明史·思宗本纪》；《明会典》相关条目；地方志、黄册/契约/河工材料 互相核对；不过，桥接条目只标示可回查的年度问题与材料链；实录有中央选择性，崇祯期尤其需要奏疏、地方、朝鲜及满文材料交叉。；本条特别核对诏令或奏议的形成与发受链。 因而这里标示的是材料能够支持的行动与制度位置，而不是对动机、地方执行或普通人经验的无保留复原。

\eventanchor{ML-1640-003}{明·1640（崇祯十三年）}
\paragraph{1640（崇祯十三年）：旱灾、蝗灾、疫病、河患与赈济的本年材料连接环境和社会压力，死亡与流离数字须谨慎处理。} 这一节点应放回本章已经展开的权力、财政、地方社会与边地关系中理解。账本把它出现的条件概括为：自然风险与地方报告促成朝廷或地方的应对记录。 随后可追踪的变化是：赈济、迁徙和损失的实际范围仍须以受灾地区材料分别核验。 可由《崇祯长编》崇祯十三年条及《明史·思宗本纪》；《明史·五行志》；受灾府县地方志与奏疏 互相核对；不过，桥接条目只标示可回查的年度问题与材料链；实录有中央选择性，崇祯期尤其需要奏疏、地方、朝鲜及满文材料交叉。 因而这里标示的是材料能够支持的行动与制度位置，而不是对动机、地方执行或普通人经验的无保留复原。

\eventanchor{ML-1640-004}{明·1640（崇祯十三年）}
\paragraph{1640（崇祯十三年）：辽饷、剿饷、练饷及地方征解的本年文书显示财政负担，定额、征收与实际流向不得混为一谈。（材料核查重点：报称信息的来源、抄传与行政分类）} 这一节点应放回本章已经展开的权力、财政、地方社会与边地关系中理解。账本把它出现的条件概括为：财政征解、交通工程或货币支付的压力使相关措施进入政务。 随后可追踪的变化是：其法定安排不等于地方执行，后续须以地域材料追踪负担与成效。 可由《崇祯长编》崇祯十三年条及《明史·思宗本纪》；《明会典》相关条目；地方志、黄册/契约/河工材料 互相核对；不过，桥接条目只标示可回查的年度问题与材料链；实录有中央选择性，崇祯期尤其需要奏疏、地方、朝鲜及满文材料交叉。；本条特别核对报称信息的来源、抄传与行政分类。 因而这里标示的是材料能够支持的行动与制度位置，而不是对动机、地方执行或普通人经验的无保留复原。

\eventanchor{ML-1640-005}{明·1640（崇祯十三年）}
\paragraph{1640（崇祯十三年）：天启末至崇祯朝的任免、奏疏和廷议构成本年中枢危机线索；崇祯无实录，必须以多类材料互校。} 这一节点应放回本章已经展开的权力、财政、地方社会与边地关系中理解。账本把它出现的条件概括为：继承、官僚协调或皇权运作的压力使问题进入朝廷程序。 随后可追踪的变化是：它影响后续人事与政策议程；动机和派系标签不能由单一叙事裁定。 可由《崇祯长编》崇祯十三年条及《明史·思宗本纪》；《明史·本纪》及相关列传；诏令、奏疏或会典版本 互相核对；不过，桥接条目只标示可回查的年度问题与材料链；实录有中央选择性，崇祯期尤其需要奏疏、地方、朝鲜及满文材料交叉。 因而这里标示的是材料能够支持的行动与制度位置，而不是对动机、地方执行或普通人经验的无保留复原。
